协同进化中的变量组同时承担两项作用：它划定交互集合，也给后端规定了必须实际求解的诱导子问题。在本文固定的 CC+CMA-ES 设定下，Q1--Q3 共同表明，仅改变分组就足以改写局部条件性；恢复隐藏的原生结构能够显著改善最终性能；显式的条件性诊断指标与终点差异之间存在稳定关联。Q4 进一步表明，结构纯度并不是终点表现的充分预测因子：在选定的病态情形下，曲率感知分组即使结构纯度更低，也仍可能在终点上优于交互优先基线 DG2。Q5 则说明，这类分组效应在更长预算下不会收敛成单一结局，而是分别表现为强烈持续、部分修复、较弱保留或退化为边界证据。综合来看，本文在特定设定下给出的经验判断是明确的：分组质量不能仅靠交互保持来刻画，诱导子问题的局部条件性是另一条与最终性能直接相关的解释维度。因此，交互保持与几何诊断承担不同职责：前者保留交互结构，后者检查分组留给后端的局部条件性，后续细化据此展开。
